\begin{theorem}[Young’s inequality]
Let $a,b\ge 0$ and let $p,q>1$ satisfy $\displaystyle\frac1p+\frac1q=1$. Then
\[
ab\;\le\;\frac{a^{p}}{p}+\frac{b^{q}}{q}\tag{1}
\]
and equality holds iff either $a=b=0$ or $a,b>0$ and $a^{p}=b^{q}$ (equivalently $a^{p-1}=b^{q-1}$).
\end{theorem}

\begin{proof}
\textbf{1. Degenerate cases.}  
If $a=0$, then $ab=0$ and the right–hand side of \eqref{1} equals $\frac{0}{p}+\frac{b^{q}}{q}\ge0$, so \eqref{1} holds.  
The same argument works when $b=0$.  
When $a=b=0$ both sides are $0$, hence equality occurs.

\medskip\noindent\textbf{2. Conjugate exponents.}  
Set
\[
\lambda:=\frac1p,\qquad \mu:=\frac1q .
\]
Because $p,q>1$ and $\frac1p+\frac1q=1$, we have
\[
0<\lambda,\mu<1,\qquad \lambda+\mu=1.\tag{2}
\]

\medskip\noindent\textbf{3. Convexity of the exponential.}  
The function $f(t)=e^{t}$ is convex on $\mathbb R$.  For any $X,Y\in\mathbb R$ and $\lambda\in[0,1]$ Jensen’s inequality (two‑point case) gives
\begin{equation}
e^{\lambda X+\mu Y}\le \lambda e^{X}+\mu e^{Y},\tag{3}
\end{equation}
with equality iff $X=Y$ (when $\lambda,\mu>0$).

\medskip\noindent\textbf{4. Choice of $X$ and $Y$.}  
Assume $a>0$ and $b>0$ (the degenerate cases have been settled).  Define
\[
X:=p\ln a,\qquad Y:=q\ln b .
\]

\medskip\noindent\textbf{5. Left–hand side of \eqref{3}.}  
Using \eqref{2},
\[
\lambda X+\mu Y
 =\frac1p(p\ln a)+\frac1q(q\ln b)
 =\ln a+\ln b
 =\ln(ab),
\]
hence
\[
e^{\lambda X+\mu Y}=e^{\ln(ab)}=ab.\tag{4}
\]

\medskip\noindent\textbf{6. Right–hand side of \eqref{3}.}  
Again by \eqref{2},
\[
\lambda e^{X}+\mu e^{Y}
 =\frac1p e^{p\ln a}+\frac1q e^{q\ln b}
 =\frac{a^{p}}{p}+\frac{b^{q}}{q}.\tag{5}
\]

\medskip\noindent\textbf{7. Conclusion of the inequality.}  
Substituting \eqref{4} and \eqref{5} into \eqref{3} yields
\[
ab\le\frac{a^{p}}{p}+\frac{b^{q}}{q},
\]
which is exactly \eqref{1}.

\medskip\noindent\textbf{8. Equality condition.}  
In \eqref{3} equality holds precisely when $X=Y$, i.e.
\[
p\ln a = q\ln b \;\Longleftrightarrow\; a^{p}=b^{q}.
\]
Together with the degenerate case $a=b=0$, the equality cases are those stated in the theorem.

Thus Young’s inequality is proved, together with its equality conditions. \qed
\end{proof}